% Fixture: statement heads in lower-case small capitals followed by a title in parentheses (no period), an exercise
% section whose items carry only their number and start past 1, a solution mark before an exercise head, and a table
% of contents listing an “Exercises” section. Compile with: tectonic exercises-smallcaps.tex
\documentclass[11pt]{book}
\usepackage{amsmath,amssymb,amsthm}
\usepackage[paperwidth=6in,paperheight=8.5in,margin=0.8in]{geometry}
\usepackage[bookmarks=true,hidelinks]{hyperref}
\newtheoremstyle{smallcaps}{12pt}{12pt}{\itshape}{}{\scshape}{}{ }{\thmname{#1}\thmnumber{ #2}\thmnote{ (#3)}}
\theoremstyle{smallcaps}
\newtheorem{definition}{definition}[chapter]
\newtheorem{lemma}[definition]{lemma}
\begin{document}
\frontmatter
\tableofcontents
\mainmatter
\chapter{Boosting}
\section{Weak learners}
A learner that does slightly better than chance is the starting point of boosting. The definition below makes this
precise; the lemma after it is used throughout the chapter.
\begin{definition}[Weak learnability]
A learning algorithm is a weak learner for a class if, for every distribution, it returns with high probability a
hypothesis whose error is at most one half minus a fixed margin.
\end{definition}
\begin{lemma}[Margin]
Every weak learner for a class of finite dimension has a margin bounded below by a constant.
\end{lemma}
\section{Boosting the weak learner}
By Definition 1.1 and Lemma 1.2, reweighting the sample after each round increases the weight of the points that the
current hypothesis misclassifies. As Exercise 1.3.5 shows, the training error then decreases exponentially fast, and
Exercise 1.3.4 asks for the number of rounds needed to reach a given error. These two exercises are the core of the
chapter; the others extend them.
\section{Exercises}
\noindent\textbf{1.3.4.} Show that the weak learner of Definition 1.1 can be boosted to any accuracy by running it
on reweighted samples, and count the rounds needed for a given error.\par\medskip
\noindent\textbf{1.3.5.} Prove that the training error of the combined hypothesis decreases exponentially fast in the
number of rounds, using Lemma 1.2 for the margin.\par\medskip
\noindent\textbf{1.3.6.} Let the margin of Lemma 1.2 shrink with the sample size. Find how many rounds are then
needed, and compare with Exercise 1.3.4.\par\medskip
\noindent\textbf{1.3.7.} Give a class with no weak learner.\par
\chapter{Margins}
\section{Margin bounds}
The margin of the combined hypothesis controls its error on new data, as the following exercise asks to show; see
also Exercise 2.1 for a direct argument.

\medskip
\noindent$\triangleleft$\ \textbf{Exercise 2.1.}[2] Prove that a large margin on the sample implies a small error
on new data, with a bound that does not depend on the number of rounds.\par\medskip
The solution uses only Definition 1.1 and the bound of Lemma 1.2.
\section{Further remarks}
Margins also explain why boosting rarely overfits, a question that Exercise 2.1 settles in the simplest case.
\end{document}
